\subsection{Attack-only}

The attack-only setting is a setting that's mainly used defensively, with a good forum on it at this \href{https://www.warzone.com/Forum/4440-usefulness-attacktransfer-orders?Offset=19}{post}. The general idea of course, is that I want to send some troops from territory A to territory B only in case my opponent has already attacked me. This is usually a positional move, though sometimes can be used to recapture a territory.
\\
\\
For example on recapturing a territory. \href{https://www.warzone.com/MultiPlayer?GameID=41541364}{Turn 4} in Partha. Goal is to outdelay the opponent then Attack-only the other territory of a double border. Extremely simple and effective idea. On the same game on Turn 7 you can also find last moves of attack-only taps, to reduce an opponent stack. Another example in this \href{https://www.warzone.com/MultiPlayer?GameID=40011579}{French Brawl game} on Turn 4. Another example in this Greece LD \href{https://www.warzone.com/MultiPlayer?GameID=39911346}{game} on Turn 14 on my very last move, as I don't want to reposition my stack due to how strong the choke-point I control is, so an attack-only is due just in case I get captured. You'll find that attack-only shines in LD templates, due to the flexibility of not rotating your armies to a weaker territory if the move fails. Another example in this \href{https://www.warzone.com/MultiPlayer?GameID=36677819}{Greece LD} game where Rufus on Turn 5 does an extra delay into an attack only move of 3 armies NE with the purpose of having two territories by the end of the turn, and give him options so as to not get eliminated.
\\
\\
Attack-only has a final very niche purpose in WR games when you want to attack a neutral to complete bonus, like in this \href{https://www.warzone.com/multiplayer?GameID=7926274}{game} as explained in the \href{https://www.warzone.com/Forum/149405-trans-attc-percentage-cards?Offset=0}{Forum}.